% GENERATED FILE. Edit canonical lessons, then run npm run generate:books.
% canonical-source-hash: fnv1a64:5a03dc997c6038fe
% canonical-lessons: RU-C358-vakansiya, RU-C358-sobesedovanie, RU-C358-kandidat, RU-C358-nanimat, RU-C358-stazh

\chapter{Looking for a job}
\label{ch:ru-looking-for-a-job}
\hlchaptermodality{358}

\begin{chapteropening}
\textbf{By the end of this chapter:} \emph{I can look for a job: talk about a vacancy, a job interview and the candidates, say who is hiring and how long I have worked.}

Complete the last lesson of chapter 358: I can look for a job: talk about a vacancy, a job interview and the candidates, say who is hiring and how long I have worked.
\end{chapteropening}

\section[\texorpdfstring{\ru{вакансия} (vakánsiya)}{вакансия (vakánsiya)}]{\ru{вакансия} (vakánsiya) — a vacancy, a job opening}

\label{lesson:RU-C358-vakansiya}



\begin{quote}
Before the new one: say the Russian for to manage, then the Russian for a team.
\end{quote}

\subsection*{You'll want to know: \ru{вакансия}}
\textbf{\ru{вакансия}} (\emph{vakánsiya}) — \textquotedblleft{}a vacancy, a job opening\textquotedblright{}.

\begin{quote}
\ru{В} \ru{этой} \ru{компании} \ru{есть} \ru{вакансия} \ru{бухгалтера}.
\end{quote}

\begin{quote}
\emph{V étoy kompánii yest vakánsiya bukhgáltera.}
\end{quote}

\begin{quote}
This company has a vacancy for an accountant.
\end{quote}

Feminine. \emph{\ru{Открытая} \ru{вакансия}} — an open position.

\begin{cousinweb}
\textbf{\ru{вакансия}} was borrowed from French \emph{vacance}, from Latin \emph{vacantia} \textquotedblleft{}emptiness, vacancy\textquotedblright{} (EtymDB), built on \emph{vacāre} \textquotedblleft{}to be empty\textquotedblright{} (Lewis \& Short). A vacancy is an empty place.
\end{cousinweb}

\begin{grammarlens}[title={What you've built}]
One more word for telling, explaining and describing at B1.
\end{grammarlens}

\subsection*{Your turn}
\begin{itemize}
\raggedright
  \item \emph{Say it:} \emph{vakánsiya}
  \item \emph{Say it:} \emph{vakánsiya}, once more
  \item \emph{Say it:} a sentence of your own with \emph{vakánsiya}
  \item \emph{Recall:} say the Russian for to manage, then the Russian for a team, then say \emph{vakánsiya} again
\end{itemize}

\begin{tcolorbox}[breakable,colback=teal!4,colframe=teal!35!black,title={Before you move on}]
What does \ru{вакансия} mean? (\textquotedblleft{}A vacancy, a job opening\textquotedblright{}.) Say it once more.
\end{tcolorbox}
\section[\texorpdfstring{\ru{собеседование} (sobesédovaniye)}{собеседование (sobesédovaniye)}]{\ru{собеседование} (sobesédovaniye) — a job interview}

\label{lesson:RU-C358-sobesedovanie}



\begin{quote}
Before the new one: say the Russian for a team, then the Russian for a vacancy.
\end{quote}

\subsection*{You'll want to know: \ru{собеседование}}
\textbf{\ru{собеседование}} (\emph{sobesédovaniye}) — \textquotedblleft{}a job interview\textquotedblright{}.

\begin{quote}
\ru{Завтра} \ru{у} \ru{меня} \ru{собеседование} \ru{в} \ru{банке}.
\end{quote}

\begin{quote}
\emph{Závtra u menyá sobesédovaniye v bánke.}
\end{quote}

\begin{quote}
Tomorrow I have a job interview at the bank.
\end{quote}

Neuter. \emph{\ru{Пройти} \ru{собеседование}} — to get through an interview.

\begin{cousinweb}
\textbf{\ru{собеседование}} is built on \emph{\ru{собеседник}} \textquotedblleft{}the person you talk with\textquotedblright{}, from \emph{\ru{со}-} \textquotedblleft{}together\textquotedblright{} and \emph{\ru{беседа}} \textquotedblleft{}a conversation\textquotedblright{}. The origin of \emph{\ru{беседа}} is not settled.
\end{cousinweb}

\begin{grammarlens}[title={What you've built}]
One more word for telling, explaining and describing at B1.
\end{grammarlens}

\subsection*{Your turn}
\begin{itemize}
\raggedright
  \item \emph{Say it:} \emph{sobesédovaniye}
  \item \emph{Say it:} \emph{sobesédovaniye}, once more
  \item \emph{Say it:} a sentence of your own with \emph{sobesédovaniye}
  \item \emph{Recall:} say the Russian for a team, then the Russian for a vacancy, then say \emph{sobesédovaniye} again
\end{itemize}

\begin{tcolorbox}[breakable,colback=teal!4,colframe=teal!35!black,title={Before you move on}]
What does \ru{собеседование} mean? (\textquotedblleft{}A job interview\textquotedblright{}.) Say it once more.
\end{tcolorbox}
\section[\texorpdfstring{\ru{кандидат} (kandidát)}{кандидат (kandidát)}]{\ru{кандидат} (kandidát) — a candidate}

\label{lesson:RU-C358-kandidat}



\begin{quote}
Before the new one: say the Russian for a vacancy, then the Russian for a job interview.
\end{quote}

\subsection*{You'll want to know: \ru{кандидат}}
\textbf{\ru{кандидат}} (\emph{kandidát}) — \textquotedblleft{}a candidate\textquotedblright{}.

\begin{quote}
\ru{На} \ru{эту} \ru{должность} \ru{пять} \ru{кандидатов}.
\end{quote}

\begin{quote}
\emph{Na étu dólzhnost pyat kandidátov.}
\end{quote}

\begin{quote}
There are five candidates for this post.
\end{quote}

Masculine. \emph{\ru{Кандидат} \ru{наук}} is also a Russian academic degree, close to a doctorate.

\begin{cousinweb}
\textbf{\ru{кандидат}} was borrowed from Polish \emph{kandydat}, from Latin \emph{candidātus} (EtymDB) \textquotedblleft{}clothed in white\textquotedblright{}, since Roman candidates for office wore a white toga (Lewis \& Short); the root is \emph{candidus} \textquotedblleft{}white\textquotedblright{}.
\end{cousinweb}

\begin{grammarlens}[title={What you've built}]
One more word for telling, explaining and describing at B1.
\end{grammarlens}

\subsection*{Your turn}
\begin{itemize}
\raggedright
  \item \emph{Say it:} \emph{kandidát}
  \item \emph{Say it:} \emph{kandidát}, once more
  \item \emph{Say it:} a sentence of your own with \emph{kandidát}
  \item \emph{Recall:} say the Russian for a vacancy, then the Russian for a job interview, then say \emph{kandidát} again
\end{itemize}

\begin{tcolorbox}[breakable,colback=teal!4,colframe=teal!35!black,title={Before you move on}]
What does \ru{кандидат} mean? (\textquotedblleft{}A candidate\textquotedblright{}.) Say it once more.
\end{tcolorbox}
\section[\texorpdfstring{\ru{нанимать} (nanimát)}{нанимать (nanimát)}]{\ru{нанимать} (nanimát) — to hire, to take on}

\label{lesson:RU-C358-nanimat}



\begin{quote}
Before the new one: say the Russian for a job interview, then the Russian for a candidate.
\end{quote}

\subsection*{You'll want to know: \ru{нанимать}}
\textbf{\ru{нанимать}} (\emph{nanimát}) — \textquotedblleft{}to hire, to take on\textquotedblright{}.

\begin{quote}
\ru{Фирма} \ru{нанимает} \ru{новых} \ru{сотрудников}.
\end{quote}

\begin{quote}
\emph{Fírma nanimáyet nóvykh sotrúdnikov.}
\end{quote}

\begin{quote}
The firm is hiring new employees.
\end{quote}

Imperfective; perfective \emph{\ru{нанять}} (\emph{\ru{найму}, \ru{наймёшь}}). An older use is to rent: \emph{\ru{нанять} \ru{квартиру}}.

\begin{cousinweb}
\textbf{\ru{нанимать}} is built from \emph{\ru{на}-} and \emph{-\ru{нимать}}, from an old verb \textquotedblleft{}to take\textquotedblright{} that also stands in \emph{\ru{снимать}} and \emph{\ru{занимать}}. To hire someone is to take them on.
\end{cousinweb}

\begin{grammarlens}[title={What you've built}]
One more word for telling, explaining and describing at B1.
\end{grammarlens}

\subsection*{Your turn}
\begin{itemize}
\raggedright
  \item \emph{Say it:} \emph{nanimát}
  \item \emph{Say it:} \emph{nanimát}, once more
  \item \emph{Say it:} a sentence of your own with \emph{nanimát}
  \item \emph{Recall:} say the Russian for a job interview, then the Russian for a candidate, then say \emph{nanimát} again
\end{itemize}

\begin{tcolorbox}[breakable,colback=teal!4,colframe=teal!35!black,title={Before you move on}]
What does \ru{нанимать} mean? (\textquotedblleft{}To hire, to take on\textquotedblright{}.) Say it once more.
\end{tcolorbox}
\section[\texorpdfstring{\ru{стаж} (stazh)}{стаж (stazh)}]{\ru{стаж} (stazh) — length of service, work experience}

\label{lesson:RU-C358-stazh}



\begin{quote}
Before the new one: say the Russian for a candidate, then the Russian for to hire.
\end{quote}

\subsection*{You'll want to know: \ru{стаж}}
\textbf{\ru{стаж}} (\emph{stazh}) — \textquotedblleft{}length of service, work experience\textquotedblright{}.

\begin{quote}
\ru{У} \ru{него} \ru{большой} \ru{стаж} \ru{работы}.
\end{quote}

\begin{quote}
\emph{U negó bolshóy stazh rabóty.}
\end{quote}

\begin{quote}
He has many years of work experience.
\end{quote}

Masculine. \emph{\ru{Стаж} \ru{работы}} — the years someone has worked, counted for a job or a pension.

\begin{cousinweb}
\textbf{\ru{стаж}} was borrowed from French \emph{stage} \textquotedblleft{}a period of training\textquotedblright{}, from Medieval Latin \emph{stagium}, from Latin \emph{stāre} \textquotedblleft{}to stand\textquotedblright{} (EtymDB).
\end{cousinweb}

\begin{grammarlens}[title={What you've built}]
One more word for telling, explaining and describing at B1.
\end{grammarlens}

\subsection*{Your turn}
\begin{itemize}
\raggedright
  \item \emph{Say it:} \emph{stazh}
  \item \emph{Say it:} \emph{stazh}, once more
  \item \emph{Say it:} a sentence of your own with \emph{stazh}
  \item \emph{Recall:} say the Russian for a candidate, then the Russian for to hire, then say \emph{stazh} again
\end{itemize}

\begin{tcolorbox}[breakable,colback=teal!4,colframe=teal!35!black,title={Before you move on}]
What does \ru{стаж} mean? (\textquotedblleft{}Length of service, work experience\textquotedblright{}.) Say it once more.
\end{tcolorbox}
